% ═══════════════════════════════════════════════════════════════════════════
% MATERIAL-09: Mystery-Karten für Speed-Dating
% 20 Karten mit Gesprächsimpulsen
% ═══════════════════════════════════════════════════════════════════════════

\documentclass[11pt, a4paper]{article}

\usepackage[utf8]{inputenc}
\usepackage[T1]{fontenc}
\usepackage[ngerman]{babel}
\usepackage{helvet}
\renewcommand{\familydefault}{\sfdefault}

\usepackage[margin=1cm]{geometry}
\usepackage[table]{xcolor}
\usepackage{tikz}
\usepackage{array}

% Farben
\definecolor{spanisch}{HTML}{c41e3a}
\definecolor{mystery}{HTML}{6a0dad}
\definecolor{hellgrau}{HTML}{f5f5f5}

\pagestyle{empty}

% Karten-Befehl
\newcommand{\mysterycard}[2]{%
\begin{tikzpicture}
\draw[mystery, line width=1.5pt, rounded corners=3pt] (0,0) rectangle (7,4);
\fill[mystery] (0,3.3) rectangle (7,4);
\node[white, font=\small\bfseries] at (3.5,3.65) {MYSTERY CARD \##1};
\node[text width=6.5cm, align=center, font=\normalsize] at (3.5,1.6) {#2};
\end{tikzpicture}%
}

\begin{document}

\begin{center}
{\Large\bfseries\textcolor{mystery}{MYSTERY-KARTEN -- Speed-Dating}}\\[0.3cm]
{\normalsize DS 09 | Zum Ausschneiden | Jeder SuS zieht 3 Karten}
\end{center}

\vspace{0.5cm}

\begin{center}

% Reihe 1
\mysterycard{1}{\textit{Pregunta:}\\\textbf{¿Cómo es tu mejor amigo/-a?}}
\hspace{0.3cm}
\mysterycard{2}{\textit{Tarea:}\\\textbf{Describe tu familia en 3 frases.}}

\vspace{0.5cm}

% Reihe 2
\mysterycard{3}{\textit{Pregunta:}\\\textbf{¿Qué hay que ser para ser un buen estudiante?}}
\hspace{0.3cm}
\mysterycard{4}{\textit{Pregunta:}\\\textbf{¿Cómo estás hoy? ¿Por qué?}}

\vspace{0.5cm}

% Reihe 3
\mysterycard{5}{\textit{Tarea:}\\\textbf{Di 5 adjetivos positivos sobre ti.}}
\hspace{0.3cm}
\mysterycard{6}{\textit{Pregunta:}\\\textbf{¿Cómo es tu profesor/-a ideal?}}

\vspace{0.5cm}

% Reihe 4
\mysterycard{7}{\textit{Pregunta:}\\\textbf{¿Tienes hermanos? ¿Cómo son?}}
\hspace{0.3cm}
\mysterycard{8}{\textit{Pregunta:}\\\textbf{¿Qué características son importantes para trabajar en grupo?}}

\end{center}

\newpage

\begin{center}

% Reihe 5
\mysterycard{9}{\textit{Pregunta:}\\\textbf{¿Qué profesión quieres tener? ¿Por qué?}}
\hspace{0.3cm}
\mysterycard{10}{\textit{Tarea:}\\\textbf{Describe a tu madre o padre en 3 frases.}}

\vspace{0.5cm}

% Reihe 6
\mysterycard{11}{\textit{Pregunta:}\\\textbf{¿Qué te gusta hacer en tu tiempo libre?}}
\hspace{0.3cm}
\mysterycard{12}{\textit{Pregunta:}\\\textbf{¿Cómo eres cuando estás con amigos?}}

\vspace{0.5cm}

% Reihe 7
\mysterycard{13}{\textit{Tarea:}\\\textbf{Di 3 cosas que hay en tu barrio.}}
\hspace{0.3cm}
\mysterycard{14}{\textit{Pregunta:}\\\textbf{¿Qué hay que hacer para ser feliz?}}

\vspace{0.5cm}

% Reihe 8
\mysterycard{15}{\textit{Pregunta:}\\\textbf{¿Cómo es tu compañero/-a de clase ideal?}}
\hspace{0.3cm}
\mysterycard{16}{\textit{Tarea:}\\\textbf{Nombra 3 cosas que te gustan.}}

\end{center}

\newpage

\begin{center}

% Reihe 9
\mysterycard{17}{\textit{Pregunta:}\\\textbf{¿Dónde vives? ¿Cómo es tu barrio?}}
\hspace{0.3cm}
\mysterycard{18}{\textit{Pregunta:}\\\textbf{¿Qué cualidades tiene un buen amigo?}}

\vspace{0.5cm}

% Reihe 10
\mysterycard{19}{\textit{Tarea:}\\\textbf{Preséntate en 4 frases.}}
\hspace{0.3cm}
\mysterycard{20}{\textit{Pregunta:}\\\textbf{¿Cómo te sientes hoy en la clase de español?}}

\end{center}

\vspace{1cm}

\hrule
\vspace{0.5cm}

\begin{center}
{\small\color{mystery}\bfseries Anleitung}
\end{center}

\begin{itemize}
    \item Karten ausschneiden und mischen
    \item Jeder SuS zieht 3 Karten zu Beginn
    \item Pro Runde (2 Min) EINE Karte spielen
    \item Die Karte gibt den Gesprächsimpuls vor
    \item Partner antwortet, dann Rollentausch
\end{itemize}

\end{document}
